%==================================================================

\subsection{Предпучок}
\leavevmode\par
\begin{defn}[предпучок]
Пусть $X$ --- топологическое пространство. \emph{Предпучком} множеств
(групп, колец $\ldots$) на $X$ называется семейство множеств
$\{\Shv{F}(U)\}$, поставленных открытым $U\subset X$, вместе с
отображениями ограничения
\[
  \res^{V}_{U}\colon\Shv{F}(V)\longrightarrow\Shv{F}(U)
  \qquad (U\subset V)
\]
такими, что
\[
  \res^{U}_{U}=\id,\qquad
  \res^{U}_{V}\circ\res^{V}_{W}=\res^{U}_{W}
  \quad\text{при}\quad U\subset V\subset W.
\]
Для $s\in\Shv{F}(V)$ пишут $s|_{U}:=\res^{V}_{U}(s)\in\Shv{F}(U)$.
\end{defn}

Короче: категория открытых множеств $\Top_X$ (объекты --- открытые
множества, морфизмы --- включения) --- это частный случай естественной
конструкции, и предпучок --- это контравариантный функтор
\[
  \Shv{F}\colon\Top_X^{\mathrm{op}}\longrightarrow \mathrm{Sets}
\]
(для абелевых групп, колец и т.\,п. --- с значениями в соответствующей
категории). Контравариантность --- это и есть то, что $\rho$ переводит
сечение на большем открытом множестве в сечение на меньшем.

\begin{rem}[что считать сечением]
Сечение --- это не функция в традиционном смысле, а <<данные, заданные
на $U$ и уменьшающиеся при уточнении>>. Типовой пример:
$\Shv{F}(U)$ --- множество непрерывных (или голоморфных, или
полиномиальных) функций на $U$, а $\rho$ --- ограничение функции.
Вся теория пучков ровно о том, что таких данных больше, чем
<<глобальных функций>>, и разрыв между ними измеряется
несвязностью $U$.
\end{rem}

\subsection{Пучок: аксиомы (F1) и (F2)}
\leavevmode\par
\begin{defn}[пучок]
Предпучок $\Shv{F}$ называется \emph{пучком}, если для любого открытого
$U$ и любого покрытия $U=\bigcup_{i\in I}U_i$ выполняются:
\begin{bullets}
\item[\textup{(F1)}] \emph{(единственность, согласованность)} если
 $s,t\in\Shv{F}(U)$ такие, что $s|_{U_i}=t|_{U_i}$ для всех $i$, то
 $s=t$;
\item[\textup{(F2)}] \emph{(склейка)} если даны $s_i\in\Shv{F}(U_i)$,
 согласованные на пересечениях,
 \[
   s_i|_{U_i\cap U_j}=s_j|_{U_i\cap U_j}\qquad\text{для всех }i,j,
 \]
 то существует $s\in\Shv{F}(U)$ с $s|_{U_i}=s_i$; причём это $s$
 определено единственным образом.
\end{bullets}
\end{defn}

Интуиция:
\begin{bullets}
\item (F1) говорит: \emph{локальное равенство --- это глобальное
 равенство}. Если два глобальных объекта неразличимы на каждом куске
 покрытия, они совпадают. То есть <<достаточно знать данные на мелких
 кусках>>.
\item (F2) говорит: \emph{согласованные локальные данные склеиваются в
 глобальные}. Обратное к (F1): <<из кусков можно собрать целое>>.
\end{bullets}
Ни одно из этих свойств не выполняется для произвольного предпучка;
именно поэтому <<пучок>> --- отдельное (и гораздо более полезное)
понятие.

\begin{bookbox}
\booksrc{Годеман, гл.~II, \S\,1, п.~1.1, с.~129 --- формулировка аксиом
 (F1), (F2) в терминах <<семейства открытых подмножеств>>; замечание:
 из (F2) в случае пустого покрытия следует, что $\Shv{F}(\varnothing)$
 --- одноэлементное множество, поэтому в (F2) достаточно рассматривать
 пары индексов с непустым пересечением.
\par\smallskip
 Манин, \S\,2.1.1---2.1.3, с.~108---110: то же определение в виде системы
 $\{\Shv{F}(U),\res^{V}_{U}\}$ и формулировка <<пучок --- пучок, если
 согласованные сечения на покрытии склеиваются и притом единственным
 образом>>. Обозначение $\Gamma(U,\Shv{F})$ для сечений.
}
\end{bookbox}

\subsection{Аксиомы пучка как точная последовательность}
Для предпучков абелевых групп аксиомы (F1), (F2) можно записать одной
строкой --- как точность последовательности. Это пригодится в конце
лекции, когда мы будем говорить о точных последовательностях пучков.

Пусть $U=\bigcup_{i\in I}U_i$ --- фиксированное покрытие. Рассмотрим
диаграмму~\ref{dgm:exact-pre}.

\begin{diagram}{Аксиомы пучка в виде точности для покрытия
 $U=\bigcup_i U_i$ (см.~также разбор ниже).\label{dgm:exact-pre}}
\begin{tikzcd}[cd style,column sep=large]
0 \ar[r] & \Shv{F}(U) \ar[r,"f"] &
\displaystyle\prod_{i\in I}\Shv{F}(U_i)
\ar[r,shift left=5pt,"\gamma"] \ar[r,shift right=5pt,"\delta"'] &
\displaystyle\prod_{i,j\in I}\Shv{F}(U_i\cap U_j)
\end{tikzcd}
\end{diagram}

Пояснение к диаграмме~\ref{dgm:exact-pre}. Читается слева
направо: стрелка $f$ <<вклеивает>> глобальное сечение в семейство
локальных, а две стрелки $\gamma$ и $\delta$ с одного множества
на другое --- это ограничение на пересечении $U_i\cap U_j$
у $i$-й и у $j$-й компоненты семейства соответственно. Точность
(равенство $\ker(\gamma-\delta)=\operatorname{Im}f$) и есть
аксиомы пучка одной строкой: «согласованные семейства --- ровно
те, что являются ограничениями глобальных сечений», --- из
неё следуют и единственность (F1), и склейка (F2); нулевой член
слева гарантирует, что $f$ ничего не теряет. Формулы компонент
--- в замечании ниже.

\begin{rem}[что означают $f,\gamma,\delta$]\label{rem:gd}
Все отображения поэлементные.
\begin{conditions}
\item Первая стрелка $f$ <<вклеивает>> глобальное сечение в семейство
 локальных:
 \begin{equation}\label{eq:f}
  f(s)=\bigl(\,\res^{U_i}_{U}(s)\,\bigr)_{i\in I}
  =\bigl(\,s|_{U_i}\,\bigr)_{i\in I}
  \in\prod_{i\in I}\Shv{F}(U_i).
 \end{equation}
\item Дальше стоят \emph{две} стрелки с одним и тем же множеством
 значений: они берут ограничение на пересечение $U_i\cap U_j$
 соответственно у $i$-й и у $j$-й компоненты семейства
 $(s_i)_{i\in I}$:
 \begin{equation}\label{eq:gd}
  \gamma\bigl((s_i)_i\bigr)_{ij}=\res^{U_i\cap U_j}_{U_i}(s_i),
  \qquad
  \delta\bigl((s_i)_i\bigr)_{ij}=\res^{U_i\cap U_j}_{U_j}(s_j).
 \end{equation}
\item Семейство $(s_i)$ называется \emph{согласованным}, если
 $\gamma\bigl((s_i)\bigr)=\delta\bigl((s_i)\bigr)$, т.\,е.\ если
 $s_i|_{U_i\cap U_j}=s_j|_{U_i\cap U_j}$ для всех $i,j$. Иными словами,
 $\gamma$ и $\delta$ --- одно и то же действие <<взять ограничение на
 пересечение>>, отнесённое к разным компонентам; их равенство и есть
 условие согласованности из (F2).
\end{conditions}
Вместо пары стрелок $\gamma,\delta$ часто пишут одну $\gamma-\delta$:
\begin{equation}\label{eq:gmd}
 (\gamma-\delta)\bigl((s_i)\bigr)_{ij}
 =\res^{U_i\cap U_j}_{U_i}(s_i)-\res^{U_j\cap U_i}_{U_j}(s_j).
\end{equation}
\end{rem}

Теперь аксиомы:
\begin{bullets}
\item \textbf{Точность в первом месте} ($f$ --- инъекция) --- это (F1):
 из $\res^{U_i}_{U}(s)=0$ для всех $i$ следует $s=0$, а для групп ---
 из равенства нулю; обобщённо: если $s|_{U_i}=t|_{U_i}$ для всех $i$,
 то $s=t$.
\item \textbf{Точность во втором месте} ($\ker(\gamma-\delta)=
 \operatorname{Im}f$) --- это (F2): слева направо --- любое согласованное
 семейство $(s_i)$ является семейством вида~\eqref{eq:f}, т.\,е.
 склеивается в глобальное $s$ (и притом единственное); справа налево ---
 глобальное сечение даёт согласованное семейство.
\item \textbf{Точность в последнем месте} не требуется: не всякое
 семейство из $\prod_{i,j}\Shv{F}(U_i\cap U_j)$ рождается из
 семейства $(s_i)$.
\end{bullets}

\begin{bookbox}
\booksrc{Манин, \S\,2.1.4, с.~110 (эта диаграмма) и \S\,2.1.5, с.~111
 (фундаментальный пример предпучка, который не пучок --- рис.~2.2 книги,
 он же рисунок~\ref{fig:circle} ниже).}
\end{bookbox}

\subsection{Пустое открытое множество}
Из (F2) для пустого покрытия $U=\varnothing$ следует, что
$\Shv{F}(\varnothing)$ --- одноэлементное множество (в абелевой
категории: нулевой элемент). Поэтому при проверке (F2) можно
рассматривать только пары индексов с непустым пересечением: если
$U_i\cap U_j=\varnothing$, условие согласованности на нём ---
тривиально, ведь ограничение в пустое множество --- единственное
отображение в однозначное множество.

\subsection{Зачем нужны именно пучки: два примера}
\leavevmode\par
\begin{exm}[склеивание двух прямых]
Пусть $X_1\simeq X_2\simeq\mathbb{R}$, $U_1\simeq U_2\simeq
\mathbb{R}\setminus\{0\}$ (две копии прямой без нуля). Склеиваем
$X_1\sqcup X_2$ по $U_1\sqcup U_2$ с помощью отображения $f$. При
$f=\id$ получается прямая с раздвоенной точкой, при $f(x)=1/x$ ---
проективная прямая (рисунок~\ref{fig:skleivanie}). Топологии
$X_1\supset U_1$ и $X_2\supset U_2$ одинаковы в обоих случаях,
и результат определяет только функция $f$. Поэтому склеивать
нужно вместе с функциями, т.\,е. работать с пучками: голая топология
разницы между $f=\id$ и $f(x)=1/x$ не видит.
\end{exm}

\begin{risunok}{Склейка двух прямых: иллюстрация к \S\,2.1 книги Манина
 (рис.~2.1).\label{fig:skleivanie}}
\resizebox{\linewidth}{!}{\begin{tikzpicture}[>=Stealth,font=\small]
  % а) Манин, рис. 2.1а: две копии прямой, стрелки --- сшиваемые пары f = id;
  % у нулей стрелки нет: 0 из первой и 0 из второй остаются двумя точками
  \draw[fig base] (-2.1,2.5) -- (2.1,2.5);
  \draw[fig base] (-2.1,1.5) -- (2.1,1.5);
  \fill (0,2.5) circle (2.4pt);
  \fill (0,1.5) circle (2.4pt);
  \node[above] at (0,2.5) {0};
  \node[below] at (0,1.5) {0};
  \foreach \x in {-1.5,-0.9,0.9,1.5}
    \draw[<->,line width=0.7pt,draw=text] (\x,1.62) -- (\x,2.38);
  \node[below,font=\footnotesize] at (0,0.75)
    {(а) $f=\id$: прямая с раздвоенной точкой};
  % б) проективная прямая
  \begin{scope}[xshift=6.6cm]
    \draw[fig base] (0,2.0) circle (1.0);
    \fill (1.0,2.0) circle (2.4pt);
    \node[below right,font=\footnotesize] at (1.05,2.1) {$\infty$};
    \node[below,font=\footnotesize] at (0,0.15)
      {(б) $f(x)=1/x$: проективная прямая};
  \end{scope}
\end{tikzpicture}}
\end{risunok}

Пояснение к рисунку~\ref{fig:skleivanie}. Слева --- результат
склейки при $f=\id$: нулевые точки двух копий не сшиваются
(стрелками показаны сшиваемые пары), получается прямая
с раздвоенной точкой: $0_1\neq 0_2$, и разнести их нельзя ---
любые окрестности пересекаются. Справа --- та же операция при
$f(x)=1/x$: <<концы>> прямых сшиваются через $1/x$, результат ---
окружность (проективная прямая), на которой обе бывшие нулевые
точки уже разделимы. Топологически $U_1,U_2$ одинаковы в обоих
случаях, различается только $f$: поэтому одной топологии
недостаточно --- нужно склеивать вместе с функциями, то есть
работать с пучками.

\begin{exm}[предпучок, не являющийся пучком]\label{exm:circle}
$X$ --- окружность, $\mathcal{O}(U)$ --- группа $\mathbb{R}$-значных
непрерывных функций на $U$, $\underline{\mathbb{Z}}$ --- <<постоянный>>
пучок. Рассмотрим предпучок
$\Shv{P}(U)=\mathcal{O}(U)/\underline{\mathbb{Z}}(U)$.
Пусть $X=U_1\cup U_2$, причём $U_1\cap U_2=V_1\sqcup V_2$ --- два
разрывных куска. Выберем $f_1\in\mathcal{O}(U_1)$, $f_2\in
\mathcal{O}(U_2)$ так, что
\[
 \res^{V_1}_{U_1}(f_1)-\res^{V_1}_{U_2}(f_2)=0,
 \qquad
 \res^{V_2}_{U_1}(f_1)-\res^{V_2}_{U_2}(f_2)=1.
\]
Классы в $\Shv{P}(V_1)$ и $\Shv{P}(V_2)$ согласованы, но глобального
сечения нет: нельзя устранить рассогласование на $V_2$, потому что
<<целое число>> нельзя поправить непрерывной функцией. Причина ---
неодносвязность окружности: аксиома (F2) здесь нарушается.
\end{exm}

\begin{risunok}{Окружность: классы $f_1$ и $f_2$ согласованы в
 $\Shv{P}$, но глобального сечения нет (Манин, рис.~2.2).
  \label{fig:circle}}
\begin{tikzpicture}[>=Stealth,font=\small]
  \draw[fig base] (0,0) circle (1.9);
  % дуги: U_1 сверху (синяя), U_2 снизу (красная), вне окружности
  \draw[fig leafA] (1.8619,1.075)
    arc[start angle=30,end angle=150,radius=2.15];
  \draw[fig sect] (1.8619,-1.075)
    arc[start angle=-30,end angle=-150,radius=2.15];
  \node[primary,font=\small] at (0,2.6) {$U_1$};
  \node[danger,font=\small] at (0,-2.6) {$U_2$};
  \node[font=\small] at (-1.25,0) {$V_1$};
  \node[font=\small] at (1.25,0) {$V_2$};
  \node[fig note,primary] at (-1.45,2.2) {$f_1$};
  \node[fig note,danger] at (1.45,-2.2) {$f_2$};
  % скачок +1 на V_2
  \draw[->,danger,line width=1pt]
    (2.2517,1.3) arc[start angle=30,end angle=-30,radius=2.6];
  \node[fig note,danger] at (3.05,0.1) {$+1$};
\end{tikzpicture}
\end{risunok}

Пояснение к рисунку~\ref{fig:circle}. Окружность распадается на
два открытых куска: $U_1$ (верхняя дуга, синяя) и $U_2$ (нижняя,
красная); их пересечение --- два разрывных куска $V_1$ слева и
$V_2$ справа. Синяя и красная линии --- это <<значения>> функций
$f_1$ и $f_2$ на соответствующих кусках. На $V_1$ концы дуг
совпадают (разность равна $0$), а на $V_2$ расходятся на $1$
(стрелка $+1$): согласованность есть, но глобального сечения
нет, потому что <<поправить>> скачок непрерывной функцией
нельзя --- в этом и состоит нарушение (F2).

%==================================================================
